%======================================================================
\section{The filtered mapping cone}\label{sec:cone}
%======================================================================
\subsection{The corner offsets}
For $n\ge1$ let $p_n\in T^2$ be the intersection point of $L_{1/(n-1)}=\{a+(n-1)b=0\}$ (for $n=1$: $L_{1/0}$) with
$L^w_0=\{b=\tfrac12\}$; it is the corner $(0,\tfrac12)$ for $n$ odd and $(\tfrac12,\tfrac12)$ for $n$ even.

\begin{lemma}[corner action]\label{lem:corner}
With the primitives of \S\ref{sec:action} ($f_{L_{1/(n-1)}}(s)=-(n-1)s^2$ on $(a,b)=s(n-1,-1)$, base point the
origin; $f_{L^w_0}(a,\tfrac12)=-a$, base point $(0,\tfrac12)$),
\[
A(p_n):=f_{L_{1/(n-1)}}(\tilde p_n)-f_{L^w_0}(\tilde p_n)\equiv\frac{n-1}{4}\pmod{\Z},
\]
independently of the lift $\tilde p_n$.
\end{lemma}
\begin{proof}
At $b=\tfrac12$ one has $s=-\tfrac12$ and $a=-(n-1)/2$, so $f_{L_{1/(n-1)}}(\tilde p_n)=-(n-1)/4$ and
$f_{L^w_0}(\tilde p_n)=-a=(n-1)/2$.  The difference is $-3(n-1)/4\equiv(n-1)/4\pmod{\Z}$.
Changing the lift along $L_{1/(n-1)}$ by $k(-(n-1),1)$ replaces $s$ by $s+k$ and changes $f_{L_{1/(n-1)}}$ by
$-(n-1)(k^2-k)\in2\Z$; changing the lift along $L_0^w$ by $m$ changes $f_{L^w_0}$ by $-m\in\Z$.
\end{proof}
The sign of the $\tfrac14$ (i.e.\ whether the offset is $(n-1)/4$ or $-(n-1)/4$ in the energy convention) is the
first thing the numerical test of \S\ref{sec:tests} fixes.

\subsection{The Dehn-twist triangle}
In $T^2$ the Dehn twist along $\{b=\tfrac12\}$ carries $L_{1/(n-1)}$ to $L^w_{1/n}=\{a+nb=\tfrac12\}$, and by
Proposition~\ref{prop:basic}(c) the translation $a\mapsto a+\tfrac12$ identifies $\Gam\cap L^w_{1/n}$ with
$\Gam\cap L_{1/n}$ (both are the flat $\SO(3)$ connections on $Y_{1/n}$).  Seidel's exact triangle for the Dehn
twist along $L^w_0$ reads, in the Fukaya category of $T^2$,
\[
L_{1/(n-1)}\xrightarrow{\ \cdot p_n\ }L^w_0\longrightarrow L^w_{1/n}\longrightarrow L_{1/(n-1)}[1],
\]
and applying $\Hom(\tilde\Gam,-)$ gives $\Hom(\tilde\Gam,L^w_{1/n})\simeq\Cone\big(\Hom(\tilde\Gam,L_{1/(n-1)})\xrightarrow{m_2(-,p_n)}\Hom(\tilde\Gam,L^w_0)\big)$.
This is the pillowcase shadow of Floer's exact triangle for the triple $(Y_{1/(n-1)},\ (Y_0,w),\ Y_{1/n})$.

\subsection{Relation to Daemi--Lidman--Miller Eismeier}
A filtered version of Floer's triangle for $\pm1$ surgery already exists: Daemi--Lidman--Miller Eismeier
\cite{DLM} prove that a cobordism map $(W,c)_*$ has degree $-2c^2$ and level $-c^2/4-\eta(W,c)$, that the Floer
triangle map has degree $0$ and level $0$, and that in the distance-two triangle the maps have level $\le\tfrac14$.
The offset $\tfrac{n-1}{4}$ below is the same quantity seen from the pillowcase: the level of a cobordism map is
the action of the corresponding corner morphism.  What is added here is (i) the generator-level description of
all three complexes and their filtrations in terms of one curve $\Gam$, and (ii) the iterated form for $1/n$.

\subsection{The theorem}
\begin{theorem}[filtered mapping cone for $1/n$ surgery]\label{thm:cone}
Assume \ref{H:alex}--\ref{H:Y}, $n\ge1$, and that $L_{1/(n-1)}$ and $L_{1/n}$ are transverse to $\Gam$ (after collar
perturbation).  Let $F_n\colon\CI(Y_{1/(n-1)})\to\CI(Y_0,w)$ be the map of the $2$-handle cobordism $W_n$.
\begin{enumerate}[label=(\alph*)]
\item \textup{(Floer, Braam--Donaldson)} $\CI(Y_{1/n})\simeq\Cone(F_n)$ as $\Z/8$-graded complexes, the grading on
$\Cone$ being $\gr_{\CI(Y_{1/(n-1)})}+\delta_n$ on the first summand and $\gr_{\CI(Y_0,w)}$ on the second, for a
constant $\delta_n\in\Z/8$ determined by the index of $W_n$ \textup{(}computed in \S\ref{sec:tests}\textup{)}.
\item \textup{(filtration)} Give $\Cone(F_n)$ the filtration $\cs$ on $\CI(Y_0,w)$ and $\cs+\tfrac{n-1}{4}+k_n$ on
$\CI(Y_{1/(n-1)})$, with $k_n\in\Z$ the unique integer for which $F_n$ is filtered and some instanton of $W_n$ has
energy $<1$.  Then there are filtered chain maps $\Cone(F_n)\rightleftarrows\CI(Y_{1/n})$ that are homotopy
inverses through filtered homotopies.  Hence the Chern--Simons filtered homotopy type of $\CI(Y_{1/n})$, and in
particular Daemi's $\Gamma_{Y_{1/n}}$ and the Nozaki--Sato--Taniguchi $r_s$, are determined by $\CI(Y_{1/(n-1)})$,
$\CI(Y_0,w)$, $F_n$ and the offset $(n-1)/4$.
\item \textup{(generator level)} The generators of all three complexes, their $\cs$ values and gradings are the
intersection data of \S\ref{sec:action}--\ref{sec:complex}; in particular, the multiset of $\cs$ values of the
irreducible flat connections on $Y_{1/n}$ is $\{A(x)\}_{x\in\Gam\cap L_{1/(n-1)}}\cup\{A(y)+\tfrac{n-1}{4}\}_{y\in\Gam\cap L^w_0}$
modulo $\Z$, up to cancelling pairs.
\end{enumerate}
\end{theorem}
\begin{proof}[Proof sketch]
(a) is the surgery exact triangle for the three fillings of slope $1/(n-1)$, $0$ (bundle $w$) and $1/n$, which
are pairwise at distance one.  (b) The maps of the triangle detection lemma (Kronheimer--Mrowka
\cite[\S7]{KM-unknot}) count instantons on $W_n$, on the second cobordism $W'_n\colon Y_0\to Y_{1/n}$, and on the
composites with up to four handles; each shifts $\cs$ by the energy, which is nonnegative, plus a topological
constant fixed by the bundle on the cobordism.  So every map is filtered once the summands are shifted by these
constants, and the constants for the two inclusions and projections of the cone are forced to be equal by the
homotopy identities.  The constant is the action of the corner morphism in the pillowcase (Lemma~\ref{lem:corner});
its instanton-theoretic computation is the topological energy of the connection on $W_n$ restricting to the flat
connections $x$ and $y$, which, by the Kirk--Klassen gluing formula applied to $W_n=(X\times I)\cup(\text{handle})$,
is the same integral.  (c) follows from Propositions \ref{prop:gens} and \ref{prop:action} and (b).
\end{proof}
What is not proved by this argument: the \emph{integer} $k_n$ (it is determined by the topology of $W_n$ but we
have not computed it; the numerical test determines it in examples), and the differential of $\CI(Y_0,w)$ itself
(an instanton count, Conjecture~\ref{conj:pillow}).

\subsection{Iteration, $Y=S^3$, negative slopes, slope $2$}
\begin{corollary}\label{cor:iter}
For $Y=S^3$ (so $\CI(Y)=0$): $\CI(S^3_1(K))\simeq\CI(S^3_0(K),w)$ with a grading shift and no filtration shift,
and for $n\ge2$ the filtered complex $\CI(S^3_{1/n}(K))$ is filtered homotopy equivalent to a twisted complex
\[
\Big(\bigoplus_{j=1}^{n}B[\delta^{(j)}]\{\tfrac{j-1}{4}+k^{(j)}\},\ \ D=\partial_B+\sum_{i<j}D_{ij}\Big),\qquad B=\CI(S^3_0(K),w),
\]
where $D_{j,j+1}$ is the Dehn-twist endomorphism $T_j=F_{j+1}\circ(\text{the equivalence }B\simeq\CI(S^3_{1/j}))$
and the $D_{ij}$, $j>i+1$, are the higher homotopies of the iterated cones.
\end{corollary}
Negative slopes $-1/n$ follow by applying the theorem to the mirror.  For slope $\pm2$ the triple
$(1/0,\ 1/1,\ 2/1)$ is at pairwise distance one, but $Y_{2}$ has a nontrivial central flat connection
(the corner $(\tfrac12,0)$ lies on $L_2$), so the irreducible-only complex must be replaced by the
$\SO(3)$-equivariant one at that corner; the pillowcase picture predicts a single extra generator with
$\cs=0$ and a corner bigon correction.  We leave $\pm2$ as a stated problem.
